% Requires: \usepackage{booktabs}, \usepackage{graphicx}, \usepackage{threeparttable}

\section{Saturation}
In an initial observation of saturation in Figure \ref{fig:Saturation}, the model cards from Anthropic and OpenAI were scraped. It appears that both Cybench and CyberGym become saturated after a year or so.
\begin{figure}
    \centering
    \includegraphics[width=1.0\linewidth]{Saturdation.png}
    \caption{Enter Caption}
    \label{fig:Saturation}
\end{figure}

Using secondary data from Epoch AI, it is able to scale the analysis. Epoch used the Inspect framework from the UK AISI. They observed a variety of benchmarks from other fields including mathematics. For each of the benchmarks, we observe the maximum scores reported for each of the releases and then resample using a step function using a monthly grid that is forward filling. The result is observed trajectories with one function per benchmark based on estimating saturation time using functional data analysis.

Based on the features of the data, the trajectories may only be partially represented as functions. This is due to the benchmarks entering the study window at different time periods and stages. For instance, some benchmarks may initially score with higher values. Of the 35 initial trajectories observed, 7 are included in this analysis to understand saturation patterns due to the period required to reach their saturation points. The other benchmarks were saturated rather quickly, making it difficult to extract signal. Analysis of the trajectories applies a parametric interpolation between reports of benchmarks.

Several approaches using functional analysis are used to observe saturation. The first approach is parametric. For each benchmark, monthly observations fit three-parameter logistic and Gompertz functions based on trajectories. The fitted rate parameter, $k$, captures the steepness of the sigmoid, and the fitted crossing of a ceiling defines a projected saturation date for benchmarks. The second approach uses anchoring for each of the points that are linearly interpolated at the 50\% performance threshold. It represents the registered curves in a cubic B-spline basis. Functional principal component analysis is used on the subset of curves within the window. The FPCA decomposes the sample of the curves into the functional mean and orthogonal modes of variation to better determine trajectory shape across the benchmarks (Table \ref{tab:fpca}). Lastly, a nonparametric model for each benchmark is applied. A higher threshold of 70\% is used to ensure the steepest segment has been observed. Velocity of saturation is estimated across months.

\begin{table}[t]
\centering
\begin{threeparttable}
\caption{Functional principal component analysis of landmark-registered saturation trajectories.}
\label{tab:fpca}
\resizebox{\columnwidth}{!}{%
\begin{tabular}{lccl}
\toprule
Component & Variance explained & Correlation\tnote{a} & Interpretation \\
\midrule
PC1 & 71.0\% & $-0.87$ & Steepness of trajectory \\
PC2 & 23.5\% & $+0.28$ & Asymmetry about midpoint \\
\bottomrule
\end{tabular}%
}
\begin{tablenotes}
\footnotesize
\item[a] Pearson correlation between the component score and the observed gain in SOTA over the registration window.
\item Note: Trajectories are registered at the interpolated 50\% crossing and represented in a cubic B-spline basis (7 basis functions) over the interval $[-12, +9]$ months. $n = 7$ benchmarks with complete coverage of the registration window.
\end{tablenotes}
\end{threeparttable}
\end{table}

Benchmark lifespans are contracting (Table \ref{tab:cohorts}). Across the benchmarks with a projected saturation date before 2030, regressing lifespan on the benchmark's entry date yields a slope of $-5.4$ months per calendar year. Each successive year of benchmark introduction is associated with roughly five and a half months less useful life. By filtering to benchmarks that have reached reporting of 95\%, the lifespan of the models is shortened by 7.5 months per year. In cohort terms, benchmarks first scored before 2023 took a median of 4.7 years to saturate. For benchmarks from 2023 and 2024, they took approximately 2.7 years to saturate. Moreover, the post-2025 cohort is estimated to be saturated after roughly 2.3 years with recent benchmarks being saturated after 13 months. One benchmark, FictionLiveBench, reached saturation in under half a year.

\begin{table}[t]
\centering
\begin{threeparttable}
\caption{Benchmark saturation dynamics by introduction cohort.}
\label{tab:cohorts}
\resizebox{\columnwidth}{!}{%
\begin{tabular}{lccccc}
\toprule
& & \multicolumn{2}{c}{Lifespan to 95\%} & & Peak velocity \\
\cmidrule(lr){3-4} \cmidrule(lr){6-6}
Cohort & Benchmarks & Median (yr) & $n$ & Median $k$ & Median (pts/mo) \\
\midrule
2019--2022 & 14 & 4.7 & 3  & 0.52 & 5.9 \\
2023       & 6  & 2.7 & 5  & 1.84 & 7.0 \\
2024       & 22 & 2.7 & 22 & 2.68 & 10.4 \\
2025--2026 & 11 & 2.3 & 9  & --- & --- \\
\bottomrule
\end{tabular}%
}
\begin{tablenotes}
\footnotesize
\item Note: Lifespan is the interval from a benchmark's first frontier score to the date its state-of-the-art (SOTA) crosses 95\% of the maximum attainable score, taken empirically where observed and from the fitted sigmoid otherwise. Sample restricted to benchmarks with at least eight scored models and a crossing date before 2030. Peak velocity is the maximum month-over-month gain in the smoothed SOTA series, computed only for benchmarks whose frontier exceeds 70\%. Entries of --- indicate the statistic could not be computed for that cohort (no 2025--2026 benchmark has yet exceeded 70\%, and the median $k$ for that cohort would rest on a single observation).
\end{tablenotes}
\end{threeparttable}
\end{table}

The contraction in lifespans is driven by steepening progress curves, and this steepening is visible under each method. Under the parametric fits, $\log(k)$ increases by approximately 0.30 per year of benchmark introduction, suggesting that the steepness of progress curves multiplies by roughly $1.35\times$ per year. The median fitted rate rises from $k = 0.52$ for pre-2023 benchmarks to $k \approx 2.6$ for the 2023--2024 cohort. The functional analysis supports the premise of this comparison: the leading principal component of the registered trajectories explains 71\% of the variation in curve shape and correlates at $-0.87$ with within-window gain.

The association between steepness and introduction date runs in the expected direction. However, it is underpowered on the seven fully observed curves. The nonparametric peak velocity estimate observes cohort medians rising from 5.9 points per month to 10.4. The corresponding regression of log peak velocity on introduction year gives approximately $1.09\times$ per year but does not reach significance. This analysis should be understood as triangulation (Table \ref{tab:regressions}). The strongly parametric estimate is statistically significant, with the weakly parametric and nonparametric estimates providing additional signal to support the claims. Per-benchmark estimates for a representative subset of the suite are reported in Table \ref{tab:benchmarks}.

\begin{table}[t]
\centering
\begin{threeparttable}
\caption{Regression estimates of the change in saturation speed on benchmark introduction date.}
\label{tab:regressions}
\resizebox{\columnwidth}{!}{%
\begin{tabular}{llcccc}
\toprule
Outcome & Sample & Slope/yr & $r$ & $p$ & $n$ \\
\midrule
Lifespan to 95\% (mo)  & Realized \& projected & $-5.4$ & $-0.53$ & $<0.001^{***}$ & 39 \\
Lifespan to 95\% (mo)  & Realized only         & $-7.5$ & $-0.91$ & $0.011^{*}$    & 6  \\
$\log$ rate $k$        & Best sigmoid fit      & $+0.30$ ($1.35\times$) & $0.57$ & $0.0001^{***}$ & 40 \\
$\log$ peak velocity   & Frontier $>$ 70\%     & $+0.09$ ($1.09\times$) & $0.29$ & $0.15$ & 27 \\
FPCA PC1 score         & Fully observed curves & $+0.11$ & $0.42$ & $0.35$ & 7 \\
\bottomrule
\end{tabular}%
}
\begin{tablenotes}
\footnotesize
\item Note: Each row reports an ordinary least squares regression of the stated outcome on the benchmark's introduction year (continuous). Rows are ordered from strongest to weakest structural assumptions. $^{*}p<0.05$; $^{**}p<0.01$; $^{***}p<0.001$.
\end{tablenotes}
\end{threeparttable}
\end{table}

\begin{table*}[t]
\centering
\begin{threeparttable}
\caption{Saturation status of selected benchmarks.}
\label{tab:benchmarks}
\resizebox{\textwidth}{!}{%
\begin{tabular}{lccccccccc}
\toprule
Benchmark & First score & Current SOTA & 95\% crossed & 95\% projected & Lifespan (yr) & Fitted model & Fitted ceiling & Rate $k$ & Peak velocity \\
\midrule
HellaSwag            & 2019-11 & 95.3\% & 2023-03 & ---     & 3.4 & Logistic & 1.00 & 0.77 & 5.4  \\
GSM8K                & 2020-06 & 94.5\% & ---     & ---     & --- & Gompertz & 0.93 & 2.74 & 18.7 \\
MMLU                 & 2021-08 & 88.1\% & ---     & ---     & --- & Gompertz & 0.86 & 2.12 & 11.4 \\
OTIS Mock AIME       & 2023-03 & 100\%  & 2025-12 & ---     & 2.7 & Gompertz & 0.94 & 3.91 & 13.4 \\
GPQA Diamond         & 2023-03 & 94.6\% & ---     & 2026-07 & 3.3 & Logistic & 1.00 & 1.16 & 5.4  \\
MATH Level 5         & 2023-06 & 98.1\% & 2025-01 & ---     & 1.6 & Logistic & 1.00 & 2.41 & 12.5 \\
Cybench              & 2024-02 & 93.0\% & ---     & 2026-12 & 2.8 & Logistic & 1.00 & 2.50 & 13.0 \\
FrontierMath         & 2024-06 & 52.4\% & ---     & 2028-01 & 3.6 & Gompertz & 0.62 & 1.58 & ---  \\
FrontierMath Tier 4  & 2024-06 & 47.9\% & ---     & 2027-07 & 3.1 & Logistic & 0.65 & 3.28 & ---  \\
ARC-AGI-2            & 2024-07 & 92.5\% & ---     & 2026-06 & 1.9 & Logistic & 0.94 & 6.20 & 15.6 \\
CritPT               & 2024-07 & 32.3\% & ---     & 2028-03 & 3.7 & Logistic & 0.38 & 3.98 & ---  \\
ARC-AGI              & 2024-09 & 98.0\% & 2026-02 & ---     & 1.4 & Logistic & 1.00 & 2.73 & 8.8  \\
HLE                  & 2024-09 & 46.4\% & ---     & 2028-01 & 3.3 & Gompertz & 0.98 & 0.97 & ---  \\
SWE-bench Verified   & 2024-11 & 83.5\% & ---     & 2027-06 & 2.6 & Logistic & 0.78 & 5.37 & 10.4 \\
GDPval               & 2024-11 & 49.7\% & ---     & 2027-11 & 3.0 & Logistic & 0.52 & 4.40 & ---  \\
Terminal-Bench       & 2025-06 & 90.2\% & ---     & 2026-08 & 1.1 & ---      & ---  & ---  & ---  \\
ExploitBench         & 2025-10 & 41.0\% & ---     & 2028-01 & 2.3 & ---      & ---  & ---  & ---  \\
OSWorld-2            & 2026-02 & 20.6\% & ---     & 2027-12 & 1.8 & ---      & ---  & ---  & ---  \\
\bottomrule
\end{tabular}%
}
\begin{tablenotes}
\footnotesize
\item Note: ``First score'' denotes the first frontier observation in the dataset, which for benchmarks predating our data source may postdate the true release date. Fitted ceilings substantially below current plausibility (e.g., FrontierMath, GDPval, CritPT) are mid-curve fitting artifacts and should be interpreted as lower bounds rather than estimates of attainable performance; GSM8K and MMLU do not cross 95\% under a free-ceiling fit, consistent with a label-noise floor. Entries of --- indicate fewer than twelve monthly observations (fit columns) or a benchmark not yet reaching the required threshold (lifespan and peak-velocity columns). Projected dates assume the fitted sigmoid continues to hold and should be interpreted as approximate.
\end{tablenotes}
\end{threeparttable}
\end{table*}
